\subsection{Bernstein's theorem}
One of the most impressive results in set theory is the so-called Bernstein's theorem. The Bernstein 
theorem basically states that even if we might not be able to exibit a bijection between two sets, such 
bijection must exist under certain conditions. This will later be useful to study the so-called "cardinality" of 
a set. \\

\begin{theorem}[Bernstein's theorem]
\label{thm:Bernstein}
    Given two sets $A$ and $B$, if there is an injection from $f : A \rightarrow B$ and an 
    injection from $g : B \rightarrow A$, then there is a bijection between the two.
\end{theorem}
\begin{proof}
    First, define the following entities (needed for the construction of the bijection):

    $$ A_0 = A - \vv{g}(B)            $$ 
    $$ A_{n+1} = \vv{g}(\vv{f}(A_n))  $$ 
    $$ A_{\infty} = \bigcup A_i       $$ \\

    
    First of all, we need to understand that $g : B \rightarrow A$ is an injection, 
    thus the very same function can be treated as a bijection when the codomain is
    restricted to the image of the domain. We can write it in symbols with the 
    following formula: \\

    $$ g : B \rightarrow A \quad\texttt{injection} 
    \Rightarrow g : B \rightarrow \vv{g}(B) \quad\texttt{bijection} $$ \\

    Since $g$ is a bijection, we can claim that $g^{-1} : \vv{g}(B) \rightarrow B $ 
    exists and is a bijection. Thus, once again, we can write the statement in 
    symbols as the following formula: \\

    $$ g : B \rightarrow A \quad\texttt{injection} 
    \Rightarrow g^{-1} : \vv{g}(B) \rightarrow A \quad\texttt{bijection} $$ \\

    We can rewrite it subsituting $\vv{g}(B)$ with $A - A_0$, since that's the very
    definition of $A_0$:

    $$ g : B \rightarrow A \quad\texttt{injection} 
    \Rightarrow g^{-1} : (A - A_0) \rightarrow A \quad\texttt{bijection} $$ \\

    We can finally restrict the domain of $g^{-1}$ to the set $A - A_{\infty}$, which 
    is a subset of $A - A_0$, hence we can write the following statement:

    $$ g : B \rightarrow A \quad\texttt{injection} 
    \Rightarrow g^{-1} : (A - A_{\infty}) \rightarrow A \quad\texttt{injection} $$ \\

    \newpage

    We can construct the bijection $h : A \rightarrow B$ and prove that it is 
    indeed a bijection (we might not be able to use it to actually compute the 
    value of $h(x)$ since, in order to do it, we might need infinite computation
    power, but we can prove that at least it exists):

    \begin{equation*}
        h(x) = 
        \begin{cases*}
            f(x) & if \( x \in A_{\infty} \)  \\
            g^{-1}(x) & otherwise
        \end{cases*}
    \end{equation*} \\


    We can also observe that $g : B \rightarrow A - A_{\infty}$ is also a bijection since $A - A_{\infty}$ is a subset of 
    $A - A_0$, hence $h$ is also a bijection for all $x \in A - A_{\infty}$ due to its very construction. \\

    We can prove that $h$ is an injection by proving that $h(x_1) = h(x_2)$ implies $x_1 = x_2$. We have to 
    reason about two cases, both of which leads to injectivity.

    $$ \forall x_1, x_2 \in A_{\infty} [ h(x_1) = f(x_1) = f(x_2) = h(x_2) \Rightarrow x_1 = x_2 ] $$
    $$ \forall x_1, x_2 \notin A_{\infty} [ h(x_1) = g^{-1}(x_1) = g^{-1}(x_2) = h(x_2) \Rightarrow x_1 = x_2 ] $$ \\
    
    Please notice that a scenario where $x_1$ and $x_2$ are not either both in $A_{\infty}$ or both are not is actually impossible under 
    the assumption that $h(x_1) = h(x_2)$, because:

    $$ \forall x_1 \in A_{\infty},\; \exists x_2 \in A \Rightarrow [ f(x_1) = h(x_1) = h(x_2) = g^{-1}(x_2) ] $$
    $$ \therefore f(x_1) = g^{-1}(x_2) \Rightarrow  g(f(x_1)) = x_2 \Rightarrow x_2 \in A_{\infty}           $$

    We can also prove that $h$ is a suriection by proving that for every $y \in B$ there is an $x \in A$ such that $h(x) = y$. To 
    find such $x$, we can reason about the following two cases:

    $$ g(y) \notin A_{\infty} \Rightarrow h(g(y)) = g^{-1}(g(y)) = y $$
    $$ g(y) \in A_{\infty} \Rightarrow g(y) \in \vv{g}(\vv{f}(A_n)) \Rightarrow f(\zeta(\vv{f}(A_n))) = y$$ \\

    It's worth noting that $g(y) \in A_{\infty}$ implies that $g(y) 
    \in A_{\infty} - A_0$ since $A_0$ is litterally defined as
    $A - \vv{g}(B)$. Hence, we can assume $g(y) \in A_{n + 1}$ \\

    By the definition of $A_{n + 1} = \vv{g}(\vv{f}(A_{n}))$, we can say that 
    $x = g(y) \in A_{n + 1}$ implies that $y = \vv{f}(A_n)$, which means we can 
    claim that the following statement holds true:

    Please notice that we are once again assuming a choice function $\zeta$ that returns a single element of a given set exist, 
    while it's easy to prove that it exist case by case, its existence cannot be proved in general. This means that the 
    Bernstein's theorem is not provable in ZF set theory alone, but it is provable in ZFC set theory (ZFC + AC). \\

\end{proof}

\newpage
